Parameter-efficient fine-tuning adapts a pretrained model by training a small number of parameters. Adapter layers~\cite{houlsby2019parameter}, bias-only tuning~\cite{benzaken2021bitfit} and low-rank adaptation (LoRA)~\cite{hu2021lora} are the best known families, and He et al.~\cite{he2021unified} give a unified view connecting several of them. LoRA freezes a pretrained weight matrix $W$ and learns a low-rank update $\frac{\alpha}{r}BA$, which can be merged back into $W$ after training. Its appeal rests on two beliefs: that the update needed for adaptation has low intrinsic dimension~\cite{aghajanyan2020intrinsic}, so a small rank $r$ suffices, and that restricting the update also limits damage to what the model already knows. The first belief is challenged where LoRA lags full fine-tuning, which motivated variants such as DoRA~\cite{liu2024dora}; the second was tested at scale by Biderman et al.~\cite{biderman2024lora}, who report that on code and mathematics LoRA learns less than full fine-tuning but forgets less of the source domain, a trade-off in the long tradition of work on catastrophic forgetting~\cite{kirkpatrick2016overcoming}.

Large-scale evidence is hard to attribute: benchmarks are noisy, pretraining data are partly unknown, and rank, learning rate and target modules interact. This paper asks the same questions in a regime where every quantity is exact. We pretrain a tiny transformer~\cite{vaswani2017attention} on one algorithmic task (ascending sort of eight digits), adapt it to two related tasks (descending sort, and reversal), and measure adaptation accuracy, trainable parameters, and retention of the pretraining task. We test four hypotheses, registered before the main runs: \textbf{H1} LoRA at rank 4 matches full fine-tuning within two points of exact match with under 15\% of the parameters; \textbf{H2} accuracy rises with rank from 1 to 8; \textbf{H3} LoRA forgets less than full fine-tuning; \textbf{H4} tuning only the last block is worse than LoRA.

Contributions, at the scale of a small CPU study: (i) a fully specified, reproducible protocol with a task token that keeps the pretraining task queryable (although the pretrained model saw only one token, so retention turned out to be at floor); (ii) results on five seeds for six reimplemented adaptation methods plus two reference rows; (iii) a target-module ablation, a training-set-size sweep and a learning-rate sweep; and (iv) an honest account of where the picture is unclear: retention is at floor for all systems so the forgetting test is uninformative, and the rank trend on reversal is confounded by optimisation failures. We make no claim about large models.
